% =============================================================================
% QUICK GUIDE - Customizing the CS Notes Template
% =============================================================================

\documentclass[]{csnotes}

\title{Quick Guide}
\author{CS Notes Template}
\date{\today}

\university{Reference Guide}
\department{Customization}
\course{LaTeX Template}
\subject{Computer Science Notes}
\academicyear{2024/2025}
\addbibresource{bibliography.bib}

\begin{document}
\maketitlepage{}
\tableofcontents
% Three separate indexes:
\listoftheorems
\listofexamples
\listofalgos

\newpage

% =============================================================================
% COMMANDS AND ENVIRONMENTS
% =============================================================================

\section{Available Commands and Environments}

\subsection{Environments for Theoretical Content}\label{sec:theoretical_environments}

\begin{definition}{Definition Name (optional)}
Use the \code{definition} environment for formal definitions.
Color: Blue
\end{definition}

\begin{theorem}{Theorem Name (optional)}
Use the \code{theorem} environment for important theorems.
Color: Red
\end{theorem}

\begin{lemma}{Lemma Name (optional)}
Use the \code{lemma} environment for lemmas and auxiliary results.
Color: Orange
\end{lemma}

\begin{proposition}{Proposition Name (optional)}
Use the \code{proposition} environment for propositions.
Color: Purple
\end{proposition}

\begin{corollary}{Corollary Name (optional)}
Use the \code{corollary} environment for corollaries.
Color: Green
\end{corollary}

\subsection{Practical Environments}\label{sec:practical_environments}

\begin{example}{Example Name (optional)}
Use the \code{example} environment for practical examples and use cases.
Color: Teal
\end{example}

\begin{observation}{Note Title (optional)}
Use the \code{observation} environment for observations about specific topics.
Color: Yellow
\end{observation}

\begin{note}{Note Title (optional)}
Use the \code{note} environment for editor notes and deeper insights.
Color: Yellow
\end{note}

\begin{warning}{Warning Title (optional)}
Use the \code{warning} environment for warnings, common mistakes, and cautions.
Color: Red
\end{warning}

\begin{algorithmdesc}{Algorithm Name (optional)}
Use the \code{algorithmdesc} environment for algorithm descriptions.
Color: Blue-gray
\end{algorithmdesc}

% =============================================================================
% PROOF DEMONSTRATION
% =============================================================================

\section{Proofs}

\begin{theorem}[label=thm:example]{Example with Proof}
This is a theorem that requires a proof.
\end{theorem}

\begin{proof}
Write the proof here. At the end, the QED symbol \(\square\) is added automatically.
\end{proof}

% Proof with a custom title
\begin{proof}[Proof by contradiction]
You can also customize the proof title.
\end{proof}

\begin{warning}{}
  If an environment has no title, the empty braces must still be included:
  \begin{lstlisting}[language=TeX]
    \begin{proposition}{}
    ...
    \end{proposition}
  \end{lstlisting}
\end{warning}

% =============================================================================
% SPECIAL COMMANDS
% =============================================================================

\section{Special Commands}

\subsection{Text Formatting}

\begin{itemize}
    \item \keyword{Keyword} --- Highlights important terms in bold blue
    \item \code{inline code} --- Formats inline code with a gray background
    \item \textbf{Regular bold} --- Standard boldface
    \item \textit{Regular italic} --- Standard italics
\end{itemize}

\subsection{Complexity Notations}

\begin{example}{Computational Complexity}
Use the dedicated commands for asymptotic notation:
\begin{itemize}
    \item \code{\textbackslash{}BigO\{n\}} produces: \(\BigO{n}\)
    \item \code{\textbackslash{}BigOmega\{n log n\}} produces: \(\BigOmega{n \log n}\)
    \item \code{\textbackslash{}BigTheta\{n\^{}2\}} produces: \(\BigTheta{n^2}\)
\end{itemize}
\end{example}

\subsection{Warning Symbols}

Commands are available to insert attention and warning symbols in the text:
\begin{itemize}
    \item \code{\textbackslash{}warningsymbol} produces: \warningsymbol
    \item \code{\textbackslash{}infosymbol} produces: \infosymbol
\end{itemize}

\subsection{Mathematical Commands}

Commands for common numeric sets are predefined:
\begin{itemize}
    \item \code{\textbackslash{}N} produces: \(\N\)
    \item \code{\textbackslash{}Z} produces: \(\Z\)
    \item \code{\textbackslash{}Q} produces: \(\Q\)
    \item \code{\textbackslash{}R} produces: \(\R\)
    \item \code{\textbackslash{}C} produces: \(\C\)
\end{itemize}

In addition to these, the template also includes:
\begin{itemize}
    \item \code{\textbackslash{}abs\{x\}} produces: \(\abs{x}\)
    \item \code{\textbackslash{}set\{x\}[condition]} produces: \(\set{x}[condition]\)
    \item \code{\textbackslash{}st} produces: \(\st\)
\end{itemize}

\subsection{Other}

Other useful commands included are:
\begin{itemize}
    \item \code{\textbackslash{}q{text}} for typographic quotation marks: \q{text}
    \item \code{\textbackslash{}qi{text}} for italic typographic quotation marks: \qi{text}
\end{itemize}

% =============================================================================
% CODE
% =============================================================================

\section{Inserting Code}

\subsection{Inline Code}

Use the command \code{\textbackslash{}code\{...\}} for short inline snippets,
for example: \code{int x = 42;} or \code{list.append(item)}.

\subsection{Code Blocks}

\begin{example}{lstlisting Syntax}
For longer code snippets, use the \code{lstlisting} environment:
\end{example}

\begin{lstlisting}[language=Python, caption=Basic example]
def hello(name):
    """Greets a person."""
    print(f"Hello, {name}!")

hello("Mario")
\end{lstlisting}

\begin{note}{Supported Languages}
Some of the supported languages are:
\begin{itemize}
    \item Python
    \item Java
    \item C, C++
    \item JavaScript
    \item SQL
    \item Haskell
    \item Bash
    \item And many more...
\end{itemize}
\end{note}

\subsection{Listing Options}

\begin{lstlisting}[
    language=C++,
    caption=Custom options,
    label=lst:example,
    numbers=left,
    firstnumber=1,
]
#include <iostream>
using namespace std;

int main() {
    cout << "Hello World!" << endl;
    return 0;
}
\end{lstlisting}

You can reference the listing with \ref{lst:example} using labels.

% =============================================================================
% MATHEMATICS
% =============================================================================
\newpage
\section{Mathematical Formulas}

\subsection{Inline Formulas}

Use \code{\textbackslash(...\textbackslash)} for inline formulas:
\(E = mc^2\), \(x = \frac{-b \pm \sqrt{b^2 - 4ac}}{2a}\).

\subsection{Displayed Formulas}

Use \code{\textbackslash[...\textbackslash]} for centered formulas:
\[
\sum_{i=1}^{n} i = \frac{n(n+1)}{2}
\]

\subsection{Equation Environments}

\begin{equation}
    \int_0^\infty e^{-x^2} dx = \frac{\sqrt{\pi}}{2}
    \label{eq:gaussian}
\end{equation}

You can reference the equation with \eqref{eq:gaussian} using labels.

\subsection{Systems of Equations}

\begin{align}
    x + y &= 5 \\
    2x - y &= 1
\end{align}

% =============================================================================
% LISTS AND ENUMERATIONS
% =============================================================================

\section{Lists}

\subsection{Bullet Lists}

\begin{itemize}
    \item First item
    \item Second item
    \begin{itemize}
        \item Sub-item
        \item Another sub-item
    \end{itemize}
    \item Third item
\end{itemize}

\subsection{Numbered Lists}

\begin{enumerate}
    \item First step
    \item Second step
    \begin{enumerate}
        \item Sub-step a
        \item Sub-step b
    \end{enumerate}
    \item Third step
\end{enumerate}

\subsection{Description Lists}

\begin{description}
    \item[Stack] LIFO structure (Last In First Out)
    \item[Queue] FIFO structure (First In First Out)
    \item[Heap] Binary tree with ordering properties
\end{description}

% =============================================================================
% TABLES AND FIGURES
% =============================================================================

\section{Tables}

\begin{table}[h]
\centering
\begin{tabular}{|l|c|c|}
\hline
\textbf{Algorithm} & \textbf{Average Complexity} & \textbf{Worst Complexity} \\
\hline
QuickSort & \(\BigO{n \log n}\) & \(\BigO{n^2}\) \\
MergeSort & \(\BigO{n \log n}\) & \(\BigO{n \log n}\) \\
HeapSort & \(\BigO{n \log n}\) & \(\BigO{n \log n}\) \\
\hline
\end{tabular}
\caption{Comparison of sorting algorithm complexity}
\label{tab:sorting}
\end{table}

Reference to the table: see \cref{tab:sorting}.

% =============================================================================
% CROSS REFERENCES
% =============================================================================

\section{Cross References}

\subsection{How to Use Labels and References}

\begin{example}{Label System}
Add labels to sections, figures, tables, and equations:
  \begin{lstlisting}[language=TeX]
  \section{Title}
  \label{sec:title}

  \begin{figure}
      ...
      \label{fig:chart}
  \end{figure}

  \begin{equation}
      ...
      \label{eq:formula}
  \end{equation}
  \end{lstlisting}
\end{example}

\begin{warning}{}
  To add references to the custom environments in \Cref{sec:theoretical_environments} and \Cref{sec:practical_environments}, use labels as optional environment properties, for example:
  \begin{lstlisting}[language=TeX]
    \begin{theorem}[label=thm:example]{Theorem Title}
    ...
    \end{theorem}
  \end{lstlisting}
  Then you can refer to the theorem with \code{\textbackslash{}cref\{thm:example\}} to include the environment type near the reference, or with \code{\textbackslash{}ref\{thm:example\}} to get only the number.
  For example: see \ref{thm:example}, or see \cref{thm:example}, or even \Cref{thm:example}.
\end{warning}

\subsection{Reference Types}

\begin{itemize}
    \item \code{\textbackslash{}ref\{label\}} --- Element number, example: \ref{sec:practical_environments}
    \item \code{\textbackslash{}eqref\{label\}} --- Equation reference with parentheses \eqref{eq:gaussian}
    \item \code{\textbackslash{}cref\{label\}} --- Smart reference (recommended), example: \cref{tab:sorting}
    \item \code{\textbackslash{}Cref\{label\}} --- Same as cref but capitalized, example: \Cref{tab:sorting}
\end{itemize}

% =============================================================================
% BIBLIOGRAPHY
% =============================================================================

\section{Bibliography (example)}

Below is a minimal example of how to use BibLaTeX in your project.
Note: add the line \verb|\addbibresource{bibliography.bib}| in the preamble of your main.tex (not here), then create the file bibliography.bib in the project root.

\begin{lstlisting}[language=TeX]
% In the preamble of main.tex:
\addbibresource{bibliography.bib}

% In the document body (citation example):
As shown in \cite{cormen2009introduction}, sorting algorithms...
% ...more text...

% At the end of the document, before \end{document}:
\printbibliography[title={References}]
\end{lstlisting}

Example bibliography.bib file:

\begin{lstlisting}[language=TeX, caption=Example bibliography.bib file]
@book{cormen2009introduction,
  title     = {Introduction to Algorithms},
  author    = {Cormen, Thomas H. and Leiserson, Charles E. and Rivest, Ronald L. and Stein, Clifford},
  year      = {2009},
  publisher = {MIT Press},
}

@article{dijkstra1959note,
  title = {A note on two problems in connexion with graphs},
  author = {Dijkstra, Edsger W.},
  journal = {Numerische Mathematik},
  year = {1959},
  volume = {1},
  pages = {269--271},
}
\end{lstlisting}




% =============================================================================
% FILE ORGANIZATION
% =============================================================================

\section{File Organization}

\subsection{Recommended Structure}

\begin{lstlisting}[language=bash]
project/
├── main.tex          % Main file
├── csnotes.cls      % Custom class
├── _chapters/       % Folder for chapters
│   ├── 0_intro.tex
│   ├── 1_chapter1.tex
│   ├── 2_chapter2.tex
│   └── ...
├── images/          % Folder for images
│   ├── figure1.png
│   ├── figure2.jpg
│   └── ...
├── code/            % Folder for source code
│   ├── example.py
│   ├── algorithm.cpp
│   └── ...
└── bibliography.bib % Bibliography file
\end{lstlisting}

\subsection{Including External Files}

\begin{example}{Input and Include}
  \begin{lstlisting}[language=TeX]
  % In main.tex
  \section{Chapter 1}
  \input{_chapters/chapter1}

  \section{Chapter 2}
  \input{_chapters/chapter2}

  % For external code
  \lstinputlisting[language=Python]{code/example.py}
  \end{lstlisting}
\end{example}

% =============================================================================
% CUSTOMIZATION
% =============================================================================

\section{Customization}

\subsection{Changing Colors}

To change the colors of the environments, edit the file \code{csnotes.cls}
in the \q{COLOR DEFINITION} section:

\begin{lstlisting}[language=TeX]
% In csnotes.cls
\definecolor{defcolor}{RGB}{0,102,204}   % Blue for definitions
\definecolor{thmcolor}{RGB}{220,50,47}   % Red for theorems
\definecolor{excolor}{RGB}{42,161,152}   % Teal for examples
% ... etc
\end{lstlisting}

\begin{note}{RGB Colors}
You can use any RGB value (0--255 for each component) or predefined names such as
\code{red}, \code{blue}, \code{green}, etc.
\end{note}

\subsection{Changing Fonts and Margins}

\begin{lstlisting}[language=TeX]
% In csnotes.cls

% Base font size (default: 11pt)
\LoadClass[a4paper,11pt]{article}

% Margins (default: 2.5cm)
\RequirePackage[margin=2.5cm, top=3cm, bottom=3cm]{geometry}
\end{lstlisting}

\subsection{Adding New Environments}

\begin{example}{Create a Custom Environment}
To add a new environment (e.g. \q{Exercise}):
\begin{lstlisting}[language=TeX]
% In csnotes.cls, after the other environments

% Define the color
\definecolor{excolor}{RGB}{255,127,0}  % Orange

% Create the environment
\newtcolorbox{exercise}[2][]{
    theorembox=excolor,
    title={Exercise~#2},
    #1
}
% Add localization for the environment type
\crefname{exercise}{exercise}{exercises}
\Crefname{exercise}{Exercise}{Exercises}
\end{lstlisting}

Then use it like this:
\begin{lstlisting}[language=TeX]
\begin{exercise}{Title}
Exercise text...
\end{exercise}
\end{lstlisting}
\end{example}

% =============================================================================
% TIPS & TRICKS
% =============================================================================

\section{Suggestions}

\begin{note}{Best Practices}
\begin{enumerate}
    \item \textbf{Consistency}: Always use the same environment for the same type of content
    \item \textbf{Brevity}: Keep boxes relatively short to avoid page break issues
    \item \textbf{Keywords}: Use \code{\textbackslash{}keyword\{\}} to highlight important terms
    \item \textbf{Organization}: Split long documents into separate files
    \item \textbf{Comments}: Use comments (\code{\%}) for personal annotations
\end{enumerate}
\end{note}

\begin{warning}{Common Mistakes}
\begin{itemize}
    \item Do not forget to compile twice to update references
    \item Always use \code{\textbackslash(...\textbackslash)} instead of \code{\$...\$} for inline math
    \item Check that the required packages are installed
    \item Do not modify \code{csnotes.cls} unless you know exactly what you are doing
\end{itemize}
\end{warning}

% =============================================================================
% ADDITIONAL RESOURCES
% =============================================================================

\section{Useful Resources}

\subsection{Documentation}

\begin{itemize}
    \item \textbf{tcolorbox}: \url{http://mirrors.ctan.org/macros/latex/contrib/tcolorbox/tcolorbox.pdf}
    \item \textbf{listings}: \url{http://mirrors.ctan.org/macros/latex/contrib/listings/listings.pdf}
    \item \textbf{amsmath}: \url{http://mirrors.ctan.org/macros/latex/required/amsmath/amsldoc.pdf}
    \item \textbf{tikz}: \url{http://mirrors.ctan.org/graphics/pgf/base/doc/pgfmanual.pdf}
\end{itemize}

\makebackcover{}

% =============================================================================
% END
% =============================================================================

\end{document}
